\section*{Chapitre \ref{ch:mx:decomposing-the-spread} --- Décomposition de la fourchette}

\begin{solution}{exo:mx:decomposing-the-spread:1}
Le mid vaut 100.02. Demi-fourchette affichée 0.02~; demi-fourchette effective $100.03-100.02=0.01$~; amélioration de prix $100.04-100.03=0.01$ par action~; demi-fourchette
réalisée $100.03-100.025=0.005$~; impact $100.025-100.02=0.005$. La fourchette effective vaut $2\times0.01/100.02$, soit 2.0 points de base.
\end{solution}

\begin{solution}{exo:mx:decomposing-the-spread:2}
Avec $\varepsilon=-1$~: demi-fourchette effective $-(49.98-50.00)=0.02$, demi-fourchette réalisée $-(49.98-49.97)=-0.01$, impact $-(49.97-50.00)=0.03$. Le fournisseur de liquidité a acheté à 49.98
un titre valant 49.97 cinq minutes plus tard~: elle a perdu un centime par action sur la transaction, avant frais et rabais.
\end{solution}

\begin{solution}{exo:mx:decomposing-the-spread:3}
Les cotations se déplacent de $\lambda S/2=0.6\times0.02=0.012$ dans le sens de la transaction~; $(1-\lambda)S/2=0.008$ reste pour le traitement (dans la forme de base, le
0.012 représente la sélection adverse et l'inventaire ensemble).
\end{solution}

\begin{solution}{exo:mx:decomposing-the-spread:4}
Fourchette implicite $2(\phi+\theta)=0.04$~; part d'information $0.012/0.020=0.6$. Un achat après un achat~: $(\phi+\theta)-(\phi+\rho\theta)=\theta(1-\rho)=0.0084$. Un achat
après une vente~: $(\phi+\theta)+(\phi+\rho\theta)=2\phi+\theta(1+\rho)=0.0316$. Le premier achat était en partie anticipé (les signes sont autocorrélés), il porte donc moins
d'information, et le rebond du cours acheteur vers le cours vendeur ajoute $2\phi$ au second.
\end{solution}

\begin{solution}{exo:mx:decomposing-the-spread:5}
Non~: l'écart, 0.28 tick, se situe à moins de deux erreurs types de la valeur à cinq minutes. Les erreurs types sont multipliées par 2.28 entre une et cinq minutes, proche
de $\sqrt5=2.24$, comme le prédit la proposition lorsque les mouvements propres du prix efficient dominent. Une fois ceux-ci retirés, l'impact vaut 0.46 tick à une
minute, et toujours 0.46 à cinq minutes (erreurs types 0.04)~: l'information a été absorbée dans la minute et y est restée.
\end{solution}

\begin{solution}{exo:mx:decomposing-the-spread:6}
En logarithmes, $\beta=\ln(100.10/99.94)^2+\ln(100.12/99.96)^2=5.12\times10^{-6}$ et $\gamma=\ln(100.12/99.94)^2=3.24\times10^{-6}$. Alors $\alpha=0.00112$ et
$S=2(e^\alpha-1)/(1+e^\alpha)\approx\alpha$~: 11.2 points de base.
\end{solution}

\begin{solution}{exo:mx:decomposing-the-spread:7}
Transactions informées~: demi-fourchette effective 0.54 tick, impact 2.11, demi-fourchette réalisée $-1.57$. Transactions des \omterm{def:mx:why-there-is-a-spread:informed}{noise traders}~: demi-fourchette effective 0.53, impact 0.045, demi-fourchette réalisée 0.49. Les \omterm{def:mx:why-there-is-a-spread:informed}{traders informés} représentent 21~\% de
le volume~: $0.21\times(-1.57)+0.79\times0.49\approx0.06$, contre 0.056 au total. Les \omterm{def:mx:why-there-is-a-spread:informed}{noise traders} paient la marge et le profit des \omterm{def:mx:why-there-is-a-spread:informed}{traders informés}.
\end{solution}

\begin{solution}{exo:mx:decomposing-the-spread:8}
Sur cinq minutes, le mid évolue avec tout ce qui se négocie après l'exécution, y compris les ordres enfants ultérieurs de l'algorithme qui poussent dans le même sens~: cela
relève de son impact et non de l'information (les transactions du métaordre se recouvrent). La moyenne à cinq minutes est aussi bruitée (exercice~5). Mesurer à des horizons plus courts,
exclure les transactions ultérieures de l'algorithme lui-même ou comparer avec un flux apparié provenant d'autres clients avant de conclure que le flux est informé.
\end{solution}

\begin{solution}{pb:mx:decomposing-the-spread:1}
\textbf{1.} $\varepsilon(p-m_t)=\varepsilon(p-m_{t+h})+\varepsilon(m_{t+h}-m_t)$~: ajouter et retrancher $\varepsilon m_{t+h}$.
\textbf{2.} \omterm{def:mx:decomposing-the-spread:quoted}{Fourchette affichée} 1.04 tick~; demi-fourchette effective 0.535 tick.
\textbf{3.} $\E[\varepsilon_t(P^\ast_t-m_t)]$, la perte espérée du fournisseur de liquidité face à ce que savait le preneur de liquidité~: 0.449 tick, 84~\%.
\textbf{4.} 2.05 tick pour les transactions informées, 0.025 pour les \omterm{def:mx:why-there-is-a-spread:informed}{noise traders}.
\textbf{5.} 0.03, 0.08, 0.35 et 0.48 tick.
\textbf{6.} Voir la proposition~: scinder l'impact en $P^\ast_{t+h}$ et $P^\ast_t$~; le terme de martingale est d'espérance nulle et de variance $\sigma^2h$.
\textbf{7.} Bruit~: l'erreur type croît de 0.07 à 0.15 tick, comme $\sqrt h$.
\textbf{8.} 0.46 tick à partir d'une minute (0.85 de la demi-fourchette effective)~: il retranche $\varepsilon(P^\ast_{t+h}-P^\ast_t)$, ce qui exige le
prix efficient.
\textbf{9.} 0.50 et 0.056 tick~: un dixième de ce qu'elle facture.
\textbf{10.} Une minute.
\textbf{11.} Prendre la différence de $p_t=\mu_t+\phi x_t+\xi_t$ et substituer $\mu_t-\mu_{t-1}=\theta(x_t-\rho x_{t-1})+u_t$.
\textbf{12.} Huang--Stoll $\lambda=0.09$ (fourchette de 1.07 tick), MRR 0.13 ($\rho=0.33$), Glosten--Harris 0.10 pour 100 actions.
\textbf{13.} Ils supposent que les cotations absorbent l'information d'une transaction dès la transaction suivante~; ici le mid se rapproche de $P^\ast$ au fil des mises à jour des cotations sur plusieurs dizaines de
secondes.
\textbf{14.} 0.14, 0.28, 0.49, 0.56 et 0.48~: davantage de retards couvrent une plus grande part de l'ajustement, puis ajoutent du bruit d'estimation.
\textbf{15.} La règle des \omterm{def:mx:why-there-is-a-spread:informed}{traders informés} dépend de la taille de l'écart $P^\ast-m$, une non-linéarité qu'une régression sur les signes gomme en moyenne.
\textbf{16.} L'amplitude n'intègre la fourchette qu'une seule fois, et la volatilité croît avec la fenêtre~; comparer les amplitudes sur un et deux jours permet de les séparer.
\textbf{17.} Corwin--Schultz 0.6 et 2.0 points de base, Abdi--Ranaldo 0.0 et 3.3, contre 1.07.
\textbf{18.} Lorsque la fourchette représente une part visible de l'amplitude de la barre (instruments moins liquides, fourchettes plus larges), et après vérification sur une période dont les fourchettes sont connues.
\textbf{19.} \emph{Résultat nommé}~: face à une part vraie de 0.84, la méthode de la demi-fourchette réalisée à une minute donne 0.89 (0.85 net du mouvement du prix efficient),
Glosten--Harris 0.10, Huang--Stoll 0.09, MRR 0.13 et le VAR de 0.14 à 0.56 selon les retards~: l'impact se stabilise vers une minute, et au-delà
les mouvements propres du prix efficient ajoutent du bruit qui croît comme $\sqrt h$.
\textbf{20.} Surtout les \omterm{def:mx:why-there-is-a-spread:informed}{noise traders}, qui ont payé les 1.57 tick par action des \omterm{def:mx:why-there-is-a-spread:informed}{traders informés} et ont laissé 0.056 au fournisseur de liquidité.
\end{solution}

\begin{solution}{iq:mx:decomposing-the-spread:1}
Demi-fourchette effective~: $\varepsilon(p-m_t)$, ce que la transaction a payé au-dessus du mid avant celle-ci. Demi-fourchette réalisée~: $\varepsilon(p-m_{t+h})$, ce que le fournisseur de liquidité a conservé au mid ultérieur. La
différence est l'\omterm{def:mx:decomposing-the-spread:impact}{impact de la transaction sur le prix} à $h$, essentiellement de la sélection adverse lorsque $h$ est bien choisi.
\iqlookfor{La convention de signe~; l'identité~; un horizon.}
\end{solution}

\begin{solution}{iq:mx:decomposing-the-spread:2}
Assez long pour que les cotations absorbent l'information de la transaction, assez court pour que la volatilité propre du prix et les transactions ultérieures ne la noient pas. Trop court~:
l'information est comptée comme fourchette gagnée. Trop long~: bruit (erreur type croissant comme $\sqrt h$) et contamination par les transactions ultérieures, y compris le même
métaordre. Tracer l'impact en fonction de l'horizon, avec les erreurs types, puis choisir l'horizon où il se stabilise.
\iqlookfor{Les deux modes de défaillance~; un choix piloté par les données.}
\end{solution}

\begin{solution}{iq:mx:decomposing-the-spread:3}
$p_t=\mu_t+\phi x_t+\xi_t$, $\mu_t=\mu_{t-1}+\theta(x_t-\rho x_{t-1})+u_t$, donc $\Delta p_t=(\phi+\theta)x_t-(\phi+\rho\theta)x_{t-1}+\dots$~; $\rho$ se déduit de
l'autocorrélation des signes, puis les deux coefficients donnent $\theta$ et $\phi$. Conditions de moments~: l'erreur orthogonale à $x_t$ et $x_{t-1}$.
\iqlookfor{La surprise dans le flux d'ordres~; l'identification par $\rho$.}
\end{solution}

\begin{solution}{iq:mx:decomposing-the-spread:4}
À une seconde, les cotations n'ont pas encore absorbé l'information~: la demi-fourchette réalisée à une seconde compte comme profit ce que la minute suivante reprend. Regarder
les mark-out à des horizons plus longs, ainsi que les pertes liées à l'inventaire et les frais.
\iqlookfor{L'horizon~; courbe de mark-out~; sélection adverse survenant après la mesure.}
\end{solution}

\begin{solution}{iq:mx:decomposing-the-spread:5}
L'impact permanent d'une transaction sur le mid des cotations, soit la réponse impulsionnelle cumulée des variations du mid à une innovation de transaction, quel que soit le chemin. Retard
nul, car les cotations sont révisées après une transaction dans l'intervalle~; l'ordre des variables identifie le choc.
\iqlookfor{Impact permanent~; réponse impulsionnelle~; hypothèse d'ordre.}
\end{solution}

\begin{solution}{iq:mx:decomposing-the-spread:6}
\omterm{def:mx:decomposing-the-spread:lowfreq}{Estimateurs de fourchette basse fréquence}~: Corwin--Schultz à partir des plus hauts et des plus bas, Abdi--Ranaldo à partir des clôtures et des amplitudes, Roll à partir des autocovariances. Les vérifier sur une période ultérieure
où les fourchettes effectives intrajournalières sont disponibles, et sur des données simulées aux fourchettes connues~; se méfier des actions liquides, où la fourchette représente une faible
part de l'amplitude.
\iqlookfor{Un estimateur nommé~; un plan de validation~; son mode de défaillance.}
\end{solution}
